%%%%%%%%%%%%%%%%%%%%%%%%%%%%%%%%%%%%%%%%%%%%%%%%%%%%%%%%%%%%
% 5.7 LEBESGUE INTEGRATION
% The Composition of Advanced Mathematics
%%%%%%%%%%%%%%%%%%%%%%%%%%%%%%%%%%%%%%%%%%%%%%%%%%%%%%%%%%%%

\section{Lebesgue Integration}
\label{sec:lebesgue-integration}

\prerequisite{
Real analysis, measure theory, measurable spaces, measurable functions,
Lebesgue measure, null sets, almost-everywhere convergence, sequences and
series, and Riemann integration.
}

\learningobjectives{
By the end of this section, the reader should be able to:
\begin{itemize}
    \item explain the motivation for Lebesgue integration;
    \item integrate indicator and nonnegative simple functions;
    \item define the Lebesgue integral of a nonnegative measurable function;
    \item work with positive and negative parts of measurable functions;
    \item determine when a function is Lebesgue integrable;
    \item understand equality almost everywhere in integration;
    \item apply the Monotone Convergence Theorem;
    \item apply Fatou's Lemma;
    \item apply the Dominated Convergence Theorem;
    \item compare Riemann and Lebesgue integration;
    \item work with \(L^1\) and \(L^p\) spaces;
    \item apply H\"older's and Minkowski's inequalities;
    \item understand essential supremum and \(L^\infty\);
    \item work with product measures;
    \item apply Tonelli's and Fubini's theorems;
    \item understand absolute continuity of the Lebesgue integral;
    \item recognize the role of Lebesgue integration throughout
          probability, functional analysis, Fourier analysis, and
          differential equations.
\end{itemize}
}

%%%%%%%%%%%%%%%%%%%%%%%%%%%%%%%%%%%%%%%%%%%%%%%%%%%%%%%%%%%%
\subsection{Why a New Integral?}
\label{subsec:why-lebesgue-integral}
%%%%%%%%%%%%%%%%%%%%%%%%%%%%%%%%%%%%%%%%%%%%%%%%%%%%%%%%%%%%

The Riemann integral is highly effective for continuous and piecewise
continuous functions.

However, modern analysis frequently studies limits of functions that are
far less regular.

Suppose

\[
f_n\to f.
\]

One naturally asks whether

\[
\boxed{
\lim_{n\to\infty}\int f_n
=
\int\lim_{n\to\infty}f_n.
}
\]

For Riemann integration, such interchanges can require restrictive
hypotheses.

Lebesgue integration is designed to interact much more effectively with
limits and measurable functions.

%%%%%%%%%%%%%%%%%%%%%%%%%%%%%%%%%%%%%%%%%%%%%%%%%%%%%%%%%%%%
\subsection{Riemann Versus Lebesgue Philosophy}
\label{subsec:riemann-versus-lebesgue-philosophy}
%%%%%%%%%%%%%%%%%%%%%%%%%%%%%%%%%%%%%%%%%%%%%%%%%%%%%%%%%%%%

The Riemann integral divides the domain into small intervals and studies
the values of the function on those intervals.

Conceptually,

\[
\boxed{
\text{Riemann: partition the domain.}
}
\]

Lebesgue integration instead groups points according to the values taken
by the function and measures the corresponding sets.

Conceptually,

\[
\boxed{
\text{Lebesgue: measure sets on which the function has given sizes.}
}
\]

This change of viewpoint makes the theory particularly compatible with
measure and convergence.

%%%%%%%%%%%%%%%%%%%%%%%%%%%%%%%%%%%%%%%%%%%%%%%%%%%%%%%%%%%%
\subsection{Indicator Functions}
\label{subsec:indicator-functions-lebesgue}
%%%%%%%%%%%%%%%%%%%%%%%%%%%%%%%%%%%%%%%%%%%%%%%%%%%%%%%%%%%%

Let

\[
(X,\mathcal{M},\mu)
\]

be a measure space and let

\[
E\in\mathcal{M}.
\]

The indicator function of \(E\) is

\[
\boxed{
\mathbf{1}_E(x)
=
\begin{cases}
1, & x\in E,\\
0, & x\notin E.
\end{cases}
}
\]

The most elementary Lebesgue integral is defined by

\[
\boxed{
\int_X \mathbf{1}_E\,d\mu
=
\mu(E).
}
\]

Thus integration begins directly from measure.

%%%%%%%%%%%%%%%%%%%%%%%%%%%%%%%%%%%%%%%%%%%%%%%%%%%%%%%%%%%%
\subsection{Integrating Constant Multiples of Indicators}
\label{subsec:constant-indicator-integral}
%%%%%%%%%%%%%%%%%%%%%%%%%%%%%%%%%%%%%%%%%%%%%%%%%%%%%%%%%%%%

If

\[
c\geq0,
\]

then

\[
\boxed{
\int_X c\mathbf{1}_E\,d\mu
=
c\,\mu(E).
}
\]

More generally, for a measurable set \(A\),

\[
\boxed{
\int_A c\,d\mu
=
c\,\mu(A).
}
\]

This provides the basic building block for simple functions.

%%%%%%%%%%%%%%%%%%%%%%%%%%%%%%%%%%%%%%%%%%%%%%%%%%%%%%%%%%%%
\subsection{Simple Functions}
\label{subsec:simple-functions-lebesgue}
%%%%%%%%%%%%%%%%%%%%%%%%%%%%%%%%%%%%%%%%%%%%%%%%%%%%%%%%%%%%

A measurable simple function takes only finitely many values.

A nonnegative simple function may be written

\[
\boxed{
\phi
=
\sum_{k=1}^{n}
a_k\mathbf{1}_{E_k},
}
\]

where

\[
a_k\geq0
\]

and the sets \(E_k\) are measurable.

We may choose the \(E_k\) to be pairwise disjoint.

%%%%%%%%%%%%%%%%%%%%%%%%%%%%%%%%%%%%%%%%%%%%%%%%%%%%%%%%%%%%
\subsection{Integral of a Nonnegative Simple Function}
\label{subsec:integral-simple-function}
%%%%%%%%%%%%%%%%%%%%%%%%%%%%%%%%%%%%%%%%%%%%%%%%%%%%%%%%%%%%

\begin{definitionbox}{Integral of a Nonnegative Simple Function}
Let

\[
\phi
=
\sum_{k=1}^{n}
a_k\mathbf{1}_{E_k},
\]

where the \(E_k\) are pairwise disjoint measurable sets and
\(a_k\geq0\).

Define

\[
\boxed{
\int_X\phi\,d\mu
=
\sum_{k=1}^{n}
a_k\mu(E_k).
}
\]
\end{definitionbox}

If \(A\in\mathcal{M}\), then

\[
\boxed{
\int_A\phi\,d\mu
=
\int_X\phi\mathbf{1}_A\,d\mu.
}
\]

%%%%%%%%%%%%%%%%%%%%%%%%%%%%%%%%%%%%%%%%%%%%%%%%%%%%%%%%%%%%
\subsection{Worked Example: Simple Function}
\label{subsec:worked-simple-integral}
%%%%%%%%%%%%%%%%%%%%%%%%%%%%%%%%%%%%%%%%%%%%%%%%%%%%%%%%%%%%

\begin{workedexamplebox}{Integrating a Simple Function}

Let

\[
\phi
=
2\mathbf{1}_{[0,1]}
+
5\mathbf{1}_{(1,3]}
\]

on \(\R\) with Lebesgue measure \(m\).

Since

\[
m([0,1])=1
\]

and

\[
m((1,3])=2,
\]

we obtain

\[
\int_{\R}\phi\,dm
=
2(1)+5(2).
\]

Therefore

\[
\begin{answerbox}
\[
\boxed{
\int_{\R}\phi\,dm=12.
}
\]
\end{answerbox}
\end{workedexamplebox}

%%%%%%%%%%%%%%%%%%%%%%%%%%%%%%%%%%%%%%%%%%%%%%%%%%%%%%%%%%%%
\subsection{Approximating Measurable Functions}
\label{subsec:simple-approximation}
%%%%%%%%%%%%%%%%%%%%%%%%%%%%%%%%%%%%%%%%%%%%%%%%%%%%%%%%%%%%

A fundamental fact is that nonnegative measurable functions can be
approximated from below by simple functions.

\begin{theorem}[Simple Function Approximation]
If

\[
f:X\to[0,\infty]
\]

is measurable, then there exists a sequence of nonnegative measurable
simple functions \((\phi_n)\) such that

\[
\boxed{
0\leq\phi_1\leq\phi_2\leq\cdots\leq f
}
\]

and

\[
\boxed{
\phi_n(x)\to f(x)
}
\]

for every \(x\in X\).
\end{theorem}

Symbolically,

\[
\boxed{
\phi_n\uparrow f.
}
\]

%%%%%%%%%%%%%%%%%%%%%%%%%%%%%%%%%%%%%%%%%%%%%%%%%%%%%%%%%%%%
\subsection{Integral of a Nonnegative Measurable Function}
\label{subsec:nonnegative-lebesgue-integral}
%%%%%%%%%%%%%%%%%%%%%%%%%%%%%%%%%%%%%%%%%%%%%%%%%%%%%%%%%%%%

\begin{definitionbox}{Lebesgue Integral for Nonnegative Functions}
Let

\[
f:X\to[0,\infty]
\]

be measurable.

Define

\[
\boxed{
\int_X f\,d\mu
=
\sup
\left\{
\int_X\phi\,d\mu:
0\leq\phi\leq f,\;
\phi\text{ simple}
\right\}.
}
\]
\end{definitionbox}

The value may be infinite.

Thus

\[
\boxed{
0\leq
\int_X f\,d\mu
\leq\infty.
}
\]

%%%%%%%%%%%%%%%%%%%%%%%%%%%%%%%%%%%%%%%%%%%%%%%%%%%%%%%%%%%%
\subsection{Monotonicity of the Integral}
\label{subsec:monotonicity-lebesgue-integral}
%%%%%%%%%%%%%%%%%%%%%%%%%%%%%%%%%%%%%%%%%%%%%%%%%%%%%%%%%%%%

If

\[
0\leq f\leq g,
\]

then

\[
\boxed{
\int_X f\,d\mu
\leq
\int_X g\,d\mu.
}
\]

More generally, if \(f\leq g\) almost everywhere and both integrals are
well defined, the same inequality holds.

%%%%%%%%%%%%%%%%%%%%%%%%%%%%%%%%%%%%%%%%%%%%%%%%%%%%%%%%%%%%
\subsection{Linearity for Nonnegative Functions}
\label{subsec:linearity-nonnegative-integrals}
%%%%%%%%%%%%%%%%%%%%%%%%%%%%%%%%%%%%%%%%%%%%%%%%%%%%%%%%%%%%

If \(f,g\geq0\) are measurable and \(a,b\geq0\), then

\[
\boxed{
\int_X(af+bg)\,d\mu
=
a\int_Xf\,d\mu
+
b\int_Xg\,d\mu.
}
\]

The equality is interpreted in the extended nonnegative real numbers.

%%%%%%%%%%%%%%%%%%%%%%%%%%%%%%%%%%%%%%%%%%%%%%%%%%%%%%%%%%%%
\subsection{Positive and Negative Parts}
\label{subsec:positive-negative-parts}
%%%%%%%%%%%%%%%%%%%%%%%%%%%%%%%%%%%%%%%%%%%%%%%%%%%%%%%%%%%%

A real-valued measurable function \(f\) can be decomposed into its positive
and negative parts.

Define

\[
\boxed{
f^+(x)
=
\max\{f(x),0\}
}
\]

and

\[
\boxed{
f^-(x)
=
\max\{-f(x),0\}.
}
\]

Then

\[
\boxed{
f=f^+-f^-,
}
\]

and

\[
\boxed{
|f|=f^++f^-.
}
\]

Also,

\[
f^+f^-=0
\]

at every point.

%%%%%%%%%%%%%%%%%%%%%%%%%%%%%%%%%%%%%%%%%%%%%%%%%%%%%%%%%%%%
\subsection{Integral of a Signed Function}
\label{subsec:signed-lebesgue-integral}
%%%%%%%%%%%%%%%%%%%%%%%%%%%%%%%%%%%%%%%%%%%%%%%%%%%%%%%%%%%%

If at least one of

\[
\int_Xf^+\,d\mu,
\qquad
\int_Xf^-\,d\mu
\]

is finite, define

\[
\boxed{
\int_Xf\,d\mu
=
\int_Xf^+\,d\mu
-
\int_Xf^-\,d\mu.
}
\]

If both positive and negative parts have infinite integral, then the
expression would have the undefined form

\[
\infty-\infty,
\]

so the integral is not defined in this sense.

%%%%%%%%%%%%%%%%%%%%%%%%%%%%%%%%%%%%%%%%%%%%%%%%%%%%%%%%%%%%
\subsection{Lebesgue Integrable Functions}
\label{subsec:lebesgue-integrable-functions}
%%%%%%%%%%%%%%%%%%%%%%%%%%%%%%%%%%%%%%%%%%%%%%%%%%%%%%%%%%%%

\begin{definitionbox}{Lebesgue Integrable Function}
A measurable function \(f:X\to\R\) is Lebesgue integrable if

\[
\boxed{
\int_X|f|\,d\mu<\infty.
}
\]

The collection of such functions is denoted

\[
\boxed{
L^1(X,\mu).
}
\]
\end{definitionbox}

Because

\[
|f|=f^++f^-,
\]

integrability is equivalent to

\[
\boxed{
\int_Xf^+\,d\mu<\infty
\quad\text{and}\quad
\int_Xf^-\,d\mu<\infty.
}
\]

%%%%%%%%%%%%%%%%%%%%%%%%%%%%%%%%%%%%%%%%%%%%%%%%%%%%%%%%%%%%
\subsection{Basic Integral Inequality}
\label{subsec:basic-integral-inequality}
%%%%%%%%%%%%%%%%%%%%%%%%%%%%%%%%%%%%%%%%%%%%%%%%%%%%%%%%%%%%

For every integrable \(f\),

\[
\boxed{
\left|
\int_Xf\,d\mu
\right|
\leq
\int_X|f|\,d\mu.
}
\]

This is the integral form of the triangle inequality.

%%%%%%%%%%%%%%%%%%%%%%%%%%%%%%%%%%%%%%%%%%%%%%%%%%%%%%%%%%%%
\subsection{Equality Almost Everywhere}
\label{subsec:equality-ae-integration}
%%%%%%%%%%%%%%%%%%%%%%%%%%%%%%%%%%%%%%%%%%%%%%%%%%%%%%%%%%%%

If

\[
f=g
\quad\text{a.e.}
\]

and \(f\) is integrable, then \(g\) is integrable and

\[
\boxed{
\int_Xf\,d\mu
=
\int_Xg\,d\mu.
}
\]

Consequently, changing the value of an integrable function on a null set
does not change its Lebesgue integral.

%%%%%%%%%%%%%%%%%%%%%%%%%%%%%%%%%%%%%%%%%%%%%%%%%%%%%%%%%%%%
\subsection{Zero Integral of a Nonnegative Function}
\label{subsec:zero-integral-nonnegative}
%%%%%%%%%%%%%%%%%%%%%%%%%%%%%%%%%%%%%%%%%%%%%%%%%%%%%%%%%%%%

\begin{theorem}
If \(f\geq0\) is measurable, then

\[
\boxed{
\int_Xf\,d\mu=0
}
\]

if and only if

\[
\boxed{
f=0
\quad\text{a.e.}
}
\]
\end{theorem}

This result is fundamental when functions differing on null sets are
identified in \(L^p\) spaces.

%%%%%%%%%%%%%%%%%%%%%%%%%%%%%%%%%%%%%%%%%%%%%%%%%%%%%%%%%%%%
\subsection{Monotone Convergence Theorem}
\label{subsec:monotone-convergence-lebesgue}
%%%%%%%%%%%%%%%%%%%%%%%%%%%%%%%%%%%%%%%%%%%%%%%%%%%%%%%%%%%%

\begin{theorem}[Monotone Convergence Theorem]
Let \((f_n)\) be a sequence of nonnegative measurable functions satisfying

\[
f_n(x)\leq f_{n+1}(x)
\]

for every \(n\) and almost every \(x\).

If

\[
f_n(x)\to f(x)
\]

almost everywhere, then

\[
\boxed{
\int_Xf\,d\mu
=
\lim_{n\to\infty}
\int_Xf_n\,d\mu.
}
\]
\end{theorem}

This theorem is also called the Beppo Levi Monotone Convergence Theorem.

It permits the interchange

\[
\boxed{
\int\lim f_n
=
\lim\int f_n
}
\]

for increasing sequences of nonnegative measurable functions.

%%%%%%%%%%%%%%%%%%%%%%%%%%%%%%%%%%%%%%%%%%%%%%%%%%%%%%%%%%%%
\subsection{Worked Example: Monotone Convergence}
\label{subsec:worked-monotone-convergence}
%%%%%%%%%%%%%%%%%%%%%%%%%%%%%%%%%%%%%%%%%%%%%%%%%%%%%%%%%%%%

\begin{workedexamplebox}{Increasing Indicators}

Let

\[
E_1\subseteq E_2\subseteq\cdots
\]

and define

\[
f_n=\mathbf{1}_{E_n}.
\]

Then

\[
f_n\uparrow
\mathbf{1}_{\bigcup_{n=1}^{\infty}E_n}.
\]

By the Monotone Convergence Theorem,

\[
\int_X
\mathbf{1}_{\bigcup E_n}\,d\mu
=
\lim_{n\to\infty}
\int_X\mathbf{1}_{E_n}\,d\mu.
\]

Therefore

\[
\begin{answerbox}
\[
\boxed{
\mu\left(
\bigcup_{n=1}^{\infty}E_n
\right)
=
\lim_{n\to\infty}\mu(E_n).
}
\]
\end{answerbox}

Thus continuity from below for measures is reflected directly in the
Monotone Convergence Theorem.
\end{workedexamplebox}

%%%%%%%%%%%%%%%%%%%%%%%%%%%%%%%%%%%%%%%%%%%%%%%%%%%%%%%%%%%%
\subsection{Fatou's Lemma}
\label{subsec:fatous-lemma}
%%%%%%%%%%%%%%%%%%%%%%%%%%%%%%%%%%%%%%%%%%%%%%%%%%%%%%%%%%%%

\begin{theorem}[Fatou's Lemma]
Let \((f_n)\) be a sequence of nonnegative measurable functions.

Then

\[
\boxed{
\int_X
\liminf_{n\to\infty}f_n\,d\mu
\leq
\liminf_{n\to\infty}
\int_Xf_n\,d\mu.
}
\]
\end{theorem}

Fatou's Lemma provides a one-sided estimate even when the sequence is not
monotone.

It is one of the principal tools for passing to limits under integrals.

%%%%%%%%%%%%%%%%%%%%%%%%%%%%%%%%%%%%%%%%%%%%%%%%%%%%%%%%%%%%
\subsection{Reverse Fatou-Type Estimate}
\label{subsec:reverse-fatou}
%%%%%%%%%%%%%%%%%%%%%%%%%%%%%%%%%%%%%%%%%%%%%%%%%%%%%%%%%%%%

Suppose there exists an integrable function \(g\) such that

\[
f_n\leq g
\]

almost everywhere.

Applying Fatou's Lemma to

\[
g-f_n
\]

can yield

\[
\boxed{
\limsup_{n\to\infty}
\int_Xf_n\,d\mu
\leq
\int_X
\limsup_{n\to\infty}f_n\,d\mu
}
\]

under appropriate hypotheses.

This idea is central to the proof of the Dominated Convergence Theorem.

%%%%%%%%%%%%%%%%%%%%%%%%%%%%%%%%%%%%%%%%%%%%%%%%%%%%%%%%%%%%
\subsection{Dominated Convergence Theorem}
\label{subsec:dominated-convergence-theorem}
%%%%%%%%%%%%%%%%%%%%%%%%%%%%%%%%%%%%%%%%%%%%%%%%%%%%%%%%%%%%

\begin{theorem}[Dominated Convergence Theorem]
Suppose \(f_n\) are measurable and

\[
f_n(x)\to f(x)
\]

almost everywhere.

Assume there exists

\[
g\in L^1(X,\mu)
\]

such that

\[
\boxed{
|f_n(x)|\leq g(x)
}
\]

for every \(n\) and almost every \(x\).

Then \(f\in L^1(X,\mu)\) and

\[
\boxed{
\lim_{n\to\infty}
\int_Xf_n\,d\mu
=
\int_Xf\,d\mu.
}
\]

Moreover,

\[
\boxed{
\int_X|f_n-f|\,d\mu\to0.
}
\]
\end{theorem}

The dominating function \(g\) provides a single integrable bound for the
entire sequence.

%%%%%%%%%%%%%%%%%%%%%%%%%%%%%%%%%%%%%%%%%%%%%%%%%%%%%%%%%%%%
\subsection{Why Domination Matters}
\label{subsec:why-domination-matters}
%%%%%%%%%%%%%%%%%%%%%%%%%%%%%%%%%%%%%%%%%%%%%%%%%%%%%%%%%%%%

Pointwise convergence alone is not sufficient to interchange limits and
integrals.

For example, on \((0,1)\), define

\[
f_n(x)
=
n\mathbf{1}_{(0,1/n)}(x).
\]

For every fixed \(x>0\),

\[
f_n(x)\to0.
\]

However,

\[
\int_0^1f_n(x)\,dx
=
n\left(\frac1n\right)
=
1.
\]

Therefore

\[
\boxed{
\lim_{n\to\infty}\int_0^1f_n\,dx
=
1
\neq
0
=
\int_0^1\lim_{n\to\infty}f_n\,dx.
}
\]

The mass concentrates near \(0\), preventing naive interchange.

%%%%%%%%%%%%%%%%%%%%%%%%%%%%%%%%%%%%%%%%%%%%%%%%%%%%%%%%%%%%
\subsection{Worked Example: Dominated Convergence}
\label{subsec:worked-dominated-convergence}
%%%%%%%%%%%%%%%%%%%%%%%%%%%%%%%%%%%%%%%%%%%%%%%%%%%%%%%%%%%%

\begin{workedexamplebox}{Passing a Limit Through an Integral}

Consider

\[
f_n(x)
=
\frac{x}{n+x}
\]

on \([0,1]\).

For each \(x\in[0,1]\),

\[
f_n(x)\to0.
\]

Also,

\[
0\leq
\frac{x}{n+x}
\leq x
\leq1.
\]

The function

\[
g(x)=1
\]

is integrable on \([0,1]\).

Therefore the Dominated Convergence Theorem gives

\[
\lim_{n\to\infty}
\int_0^1
\frac{x}{n+x}\,dx
=
\int_0^1 0\,dx.
\]

Hence

\[
\begin{answerbox}
\[
\boxed{
\lim_{n\to\infty}
\int_0^1
\frac{x}{n+x}\,dx
=
0.
}
\]
\end{answerbox}
\end{workedexamplebox}

%%%%%%%%%%%%%%%%%%%%%%%%%%%%%%%%%%%%%%%%%%%%%%%%%%%%%%%%%%%%
\subsection{Comparison of the Three Major Convergence Results}
\label{subsec:comparison-convergence-theorems}
%%%%%%%%%%%%%%%%%%%%%%%%%%%%%%%%%%%%%%%%%%%%%%%%%%%%%%%%%%%%

The three fundamental convergence results play different roles.

\[
\boxed{
\begin{array}{c|c}
\text{Theorem} & \text{Main Hypothesis}\\
\hline
\text{Monotone Convergence} &
0\leq f_n\uparrow f\\
\text{Fatou's Lemma} &
f_n\geq0\\
\text{Dominated Convergence} &
f_n\to f\text{ a.e. and }|f_n|\leq g\in L^1
\end{array}
}
\]

The Monotone and Dominated Convergence Theorems give equality of a limit
and an integral.

Fatou's Lemma gives a one-sided inequality.

%%%%%%%%%%%%%%%%%%%%%%%%%%%%%%%%%%%%%%%%%%%%%%%%%%%%%%%%%%%%
\subsection{Riemann and Lebesgue Integrals}
\label{subsec:riemann-lebesgue-comparison}
%%%%%%%%%%%%%%%%%%%%%%%%%%%%%%%%%%%%%%%%%%%%%%%%%%%%%%%%%%%%

Lebesgue integration extends Riemann integration.

\begin{theorem}
If a bounded function

\[
f:[a,b]\to\R
\]

is Riemann integrable, then it is Lebesgue integrable and the two integrals
agree:

\[
\boxed{
\int_a^b f(x)\,dx
=
\int_{[a,b]}f\,dm.
}
\]
\end{theorem}

Thus Lebesgue integration does not replace the numerical value of the
Riemann integral where the latter already works.

It extends the class of integrable functions and improves convergence
theory.

%%%%%%%%%%%%%%%%%%%%%%%%%%%%%%%%%%%%%%%%%%%%%%%%%%%%%%%%%%%%
\subsection{Lebesgue's Criterion for Riemann Integrability}
\label{subsec:lebesgue-criterion-riemann}
%%%%%%%%%%%%%%%%%%%%%%%%%%%%%%%%%%%%%%%%%%%%%%%%%%%%%%%%%%%%

\begin{theorem}[Lebesgue's Criterion]
A bounded function

\[
f:[a,b]\to\R
\]

is Riemann integrable if and only if its set of discontinuities has
Lebesgue measure zero.
\end{theorem}

This gives a precise measure-theoretic characterization of Riemann
integrability.

For example, a bounded function with finitely or countably many
discontinuities is Riemann integrable.

%%%%%%%%%%%%%%%%%%%%%%%%%%%%%%%%%%%%%%%%%%%%%%%%%%%%%%%%%%%%
\subsection{The Dirichlet Function Revisited}
\label{subsec:dirichlet-lebesgue-integral}
%%%%%%%%%%%%%%%%%%%%%%%%%%%%%%%%%%%%%%%%%%%%%%%%%%%%%%%%%%%%

Consider

\[
f(x)
=
\mathbf{1}_{\Q}(x)
\]

on \([0,1]\).

This function is discontinuous everywhere, so it is not Riemann
integrable.

However,

\[
m(\Q\cap[0,1])=0.
\]

Therefore

\[
f=0
\quad\text{a.e.}
\]

and hence

\[
\boxed{
\int_{[0,1]}
\mathbf{1}_{\Q}\,dm
=
0.
}
\]

Thus the function is Lebesgue integrable.

%%%%%%%%%%%%%%%%%%%%%%%%%%%%%%%%%%%%%%%%%%%%%%%%%%%%%%%%%%%%
\subsection{Absolute Continuity of the Integral}
\label{subsec:absolute-continuity-integral}
%%%%%%%%%%%%%%%%%%%%%%%%%%%%%%%%%%%%%%%%%%%%%%%%%%%%%%%%%%%%

\begin{theorem}[Absolute Continuity of the Lebesgue Integral]
If

\[
f\in L^1(X,\mu),
\]

then for every \(\varepsilon>0\), there exists \(\delta>0\) such that

\[
\boxed{
\mu(E)<\delta
\Longrightarrow
\int_E|f|\,d\mu<\varepsilon
}
\]

for every measurable \(E\).
\end{theorem}

Thus an integrable function cannot place a substantial amount of integral
mass on arbitrarily small measurable sets.

%%%%%%%%%%%%%%%%%%%%%%%%%%%%%%%%%%%%%%%%%%%%%%%%%%%%%%%%%%%%
\subsection{The \(L^1\) Norm}
\label{subsec:l1-norm}
%%%%%%%%%%%%%%%%%%%%%%%%%%%%%%%%%%%%%%%%%%%%%%%%%%%%%%%%%%%%

For

\[
f\in L^1(X,\mu),
\]

define

\[
\boxed{
\|f\|_1
=
\int_X|f|\,d\mu.
}
\]

The double bars

\[
\|\cdot\|
\]

denote a \emph{norm}.

Strictly speaking, measurable functions that agree almost everywhere must
be identified in order for this to be a genuine norm.

If functions are not identified almost everywhere, then

\[
\|f\|_1=0
\]

only implies

\[
f=0
\quad\text{a.e.},
\]

rather than everywhere.

%%%%%%%%%%%%%%%%%%%%%%%%%%%%%%%%%%%%%%%%%%%%%%%%%%%%%%%%%%%%
\subsection{\(L^p\) Spaces}
\label{subsec:lp-spaces-introduction}
%%%%%%%%%%%%%%%%%%%%%%%%%%%%%%%%%%%%%%%%%%%%%%%%%%%%%%%%%%%%

For

\[
1\leq p<\infty,
\]

define

\[
\boxed{
L^p(X,\mu)
=
\left\{
f:
f\text{ measurable and }
\int_X|f|^p\,d\mu<\infty
\right\},
}
\]

where functions equal almost everywhere are identified.

The \(L^p\) norm is

\[
\boxed{
\|f\|_p
=
\left(
\int_X|f|^p\,d\mu
\right)^{1/p}.
}
\]

Important cases include

\[
\boxed{
L^1,
\qquad
L^2,
\qquad
L^\infty.
}
\]

%%%%%%%%%%%%%%%%%%%%%%%%%%%%%%%%%%%%%%%%%%%%%%%%%%%%%%%%%%%%
\subsection{The Space \(L^2\)}
\label{subsec:l2-space}
%%%%%%%%%%%%%%%%%%%%%%%%%%%%%%%%%%%%%%%%%%%%%%%%%%%%%%%%%%%%

For

\[
p=2,
\]

the norm becomes

\[
\boxed{
\|f\|_2
=
\left(
\int_X|f|^2\,d\mu
\right)^{1/2}.
}
\]

The space \(L^2\) has additional geometric structure because it arises
from the inner product

\[
\boxed{
\langle f,g\rangle
=
\int_X f\overline{g}\,d\mu
}
\]

for complex-valued functions.

For real-valued functions this reduces to

\[
\boxed{
\langle f,g\rangle
=
\int_Xfg\,d\mu.
}
\]

This makes \(L^2\) central to Fourier analysis and functional analysis.

%%%%%%%%%%%%%%%%%%%%%%%%%%%%%%%%%%%%%%%%%%%%%%%%%%%%%%%%%%%%
\subsection{Essential Supremum}
\label{subsec:essential-supremum}
%%%%%%%%%%%%%%%%%%%%%%%%%%%%%%%%%%%%%%%%%%%%%%%%%%%%%%%%%%%%

Ordinary supremum is sensitive to values attained even at a single point.

Measure theory instead often uses the essential supremum.

\begin{definitionbox}{Essential Supremum}
For a measurable function \(f\), the essential supremum of \(|f|\) is

\[
\boxed{
\operatorname*{ess\,sup}_{x\in X}|f(x)|
=
\inf
\{
M\geq0:
|f(x)|\leq M
\text{ a.e.}
\}.
}
\]
\end{definitionbox}

Values occurring only on null sets do not affect the essential supremum.

%%%%%%%%%%%%%%%%%%%%%%%%%%%%%%%%%%%%%%%%%%%%%%%%%%%%%%%%%%%%
\subsection{The Space \(L^\infty\)}
\label{subsec:l-infinity}
%%%%%%%%%%%%%%%%%%%%%%%%%%%%%%%%%%%%%%%%%%%%%%%%%%%%%%%%%%%%

A measurable function belongs to \(L^\infty\) if it is essentially
bounded.

Its norm is

\[
\boxed{
\|f\|_\infty
=
\operatorname*{ess\,sup}|f|.
}
\]

For example, if

\[
f(x)=
\begin{cases}
1000, & x=0,\\
1, & x\neq0
\end{cases}
\]

on \([0,1]\), then

\[
\boxed{
\|f\|_\infty=1
}
\]

with respect to Lebesgue measure.

The exceptional value at one point is ignored because that point has
measure zero.

%%%%%%%%%%%%%%%%%%%%%%%%%%%%%%%%%%%%%%%%%%%%%%%%%%%%%%%%%%%%
\subsection{H\"older's Inequality}
\label{subsec:holders-inequality}
%%%%%%%%%%%%%%%%%%%%%%%%%%%%%%%%%%%%%%%%%%%%%%%%%%%%%%%%%%%%

Let

\[
1<p,q<\infty
\]

satisfy

\[
\boxed{
\frac1p+\frac1q=1.
}
\]

The numbers \(p\) and \(q\) are called conjugate exponents.

\begin{theorem}[H\"older's Inequality]
If

\[
f\in L^p
\]

and

\[
g\in L^q,
\]

then

\[
fg\in L^1
\]

and

\[
\boxed{
\int_X|fg|\,d\mu
\leq
\|f\|_p\|g\|_q.
}
\]
\end{theorem}

For

\[
p=q=2,
\]

H\"older's inequality becomes the Cauchy--Schwarz inequality:

\[
\boxed{
\left|
\int_Xf\overline{g}\,d\mu
\right|
\leq
\|f\|_2\|g\|_2.
}
\]

%%%%%%%%%%%%%%%%%%%%%%%%%%%%%%%%%%%%%%%%%%%%%%%%%%%%%%%%%%%%
\subsection{Minkowski's Inequality}
\label{subsec:minkowski-inequality}
%%%%%%%%%%%%%%%%%%%%%%%%%%%%%%%%%%%%%%%%%%%%%%%%%%%%%%%%%%%%

\begin{theorem}[Minkowski's Inequality]
For

\[
1\leq p\leq\infty,
\]

\[
\boxed{
\|f+g\|_p
\leq
\|f\|_p+\|g\|_p.
}
\]
\end{theorem}

This is the triangle inequality for the \(L^p\) norm.

Together with H\"older's inequality, it establishes the basic normed-space
structure of \(L^p\).

%%%%%%%%%%%%%%%%%%%%%%%%%%%%%%%%%%%%%%%%%%%%%%%%%%%%%%%%%%%%
\subsection{Completeness of \(L^p\)}
\label{subsec:completeness-lp}
%%%%%%%%%%%%%%%%%%%%%%%%%%%%%%%%%%%%%%%%%%%%%%%%%%%%%%%%%%%%

\begin{theorem}[Riesz--Fischer Completeness]
For

\[
1\leq p\leq\infty,
\]

the space

\[
\boxed{
L^p(X,\mu)
}
\]

is complete with respect to the \(L^p\) norm.
\end{theorem}

Therefore every \(L^p\)-Cauchy sequence converges in \(L^p\) to an element
of \(L^p\).

A complete normed vector space is called a Banach space.

Thus

\[
\boxed{
L^p
\text{ is a Banach space}.
}
\]

The special space \(L^2\) is a Hilbert space.

%%%%%%%%%%%%%%%%%%%%%%%%%%%%%%%%%%%%%%%%%%%%%%%%%%%%%%%%%%%%
\subsection{Convergence in \(L^p\)}
\label{subsec:lp-convergence}
%%%%%%%%%%%%%%%%%%%%%%%%%%%%%%%%%%%%%%%%%%%%%%%%%%%%%%%%%%%%

A sequence \((f_n)\) converges to \(f\) in \(L^p\) if

\[
\boxed{
\|f_n-f\|_p\to0.
}
\]

For \(1\leq p<\infty\), this means

\[
\boxed{
\int_X|f_n-f|^p\,d\mu\to0.
}
\]

This is another important mode of convergence, distinct from pointwise,
uniform, almost-everywhere, and in-measure convergence.

%%%%%%%%%%%%%%%%%%%%%%%%%%%%%%%%%%%%%%%%%%%%%%%%%%%%%%%%%%%%
\subsection{\(L^p\) Convergence and Convergence in Measure}
\label{subsec:lp-to-measure}
%%%%%%%%%%%%%%%%%%%%%%%%%%%%%%%%%%%%%%%%%%%%%%%%%%%%%%%%%%%%

For

\[
1\leq p<\infty,
\]

if

\[
f_n\to f
\quad\text{in }L^p,
\]

then

\[
\boxed{
f_n\to f
\quad\text{in measure}.
}
\]

This follows from Markov-type estimates such as

\[
\boxed{
\mu(
\{|f_n-f|>\varepsilon\}
)
\leq
\frac{\|f_n-f\|_p^p}{\varepsilon^p}.
}
\]

Thus \(L^p\) convergence provides quantitative control over convergence in
measure.

%%%%%%%%%%%%%%%%%%%%%%%%%%%%%%%%%%%%%%%%%%%%%%%%%%%%%%%%%%%%
\subsection{Almost-Everywhere Convergence and \(L^p\)}
\label{subsec:ae-versus-lp}
%%%%%%%%%%%%%%%%%%%%%%%%%%%%%%%%%%%%%%%%%%%%%%%%%%%%%%%%%%%%

Almost-everywhere convergence alone does not imply \(L^p\) convergence.

For example,

\[
f_n(x)
=
n\mathbf{1}_{(0,1/n)}(x)
\]

satisfies

\[
f_n(x)\to0
\]

for every \(x>0\), but

\[
\|f_n\|_1=1.
\]

Hence

\[
f_n\not\to0
\]

in \(L^1\).

Additional control, such as domination, is required.

%%%%%%%%%%%%%%%%%%%%%%%%%%%%%%%%%%%%%%%%%%%%%%%%%%%%%%%%%%%%
\subsection{Product Measure Spaces}
\label{subsec:product-measure-spaces}
%%%%%%%%%%%%%%%%%%%%%%%%%%%%%%%%%%%%%%%%%%%%%%%%%%%%%%%%%%%%

Let

\[
(X,\mathcal{M},\mu)
\]

and

\[
(Y,\mathcal{N},\nu)
\]

be sigma-finite measure spaces.

The product sigma-algebra is

\[
\boxed{
\mathcal{M}\otimes\mathcal{N},
}
\]

generated by measurable rectangles

\[
A\times B,
\qquad
A\in\mathcal{M},
\quad
B\in\mathcal{N}.
\]

The product measure

\[
\mu\times\nu
\]

satisfies

\[
\boxed{
(\mu\times\nu)(A\times B)
=
\mu(A)\nu(B).
}
\]

%%%%%%%%%%%%%%%%%%%%%%%%%%%%%%%%%%%%%%%%%%%%%%%%%%%%%%%%%%%%
\subsection{Sections of Functions}
\label{subsec:function-sections}
%%%%%%%%%%%%%%%%%%%%%%%%%%%%%%%%%%%%%%%%%%%%%%%%%%%%%%%%%%%%

For a function

\[
f:X\times Y\to\overline{\R},
\]

fixing one variable produces a section.

For fixed \(x\in X\),

\[
\boxed{
f_x(y)=f(x,y).
}
\]

For fixed \(y\in Y\),

\[
\boxed{
f^y(x)=f(x,y).
}
\]

Iterated integrals integrate these sections one variable at a time.

%%%%%%%%%%%%%%%%%%%%%%%%%%%%%%%%%%%%%%%%%%%%%%%%%%%%%%%%%%%%
\subsection{Tonelli's Theorem}
\label{subsec:tonellis-theorem}
%%%%%%%%%%%%%%%%%%%%%%%%%%%%%%%%%%%%%%%%%%%%%%%%%%%%%%%%%%%%

\begin{theorem}[Tonelli's Theorem]
Let

\[
f:X\times Y\to[0,\infty]
\]

be measurable.

Then

\[
\boxed{
\int_{X\times Y}
f\,d(\mu\times\nu)
=
\int_X
\left(
\int_Yf(x,y)\,d\nu(y)
\right)
d\mu(x)
}
\]

and

\[
\boxed{
\int_{X\times Y}
f\,d(\mu\times\nu)
=
\int_Y
\left(
\int_Xf(x,y)\,d\mu(x)
\right)
d\nu(y).
}
\]

The values are permitted to be infinite.
\end{theorem}

Tonelli's theorem is especially powerful because nonnegativity is enough;
integrability need not be known beforehand.

%%%%%%%%%%%%%%%%%%%%%%%%%%%%%%%%%%%%%%%%%%%%%%%%%%%%%%%%%%%%
\subsection{Fubini's Theorem}
\label{subsec:fubinis-theorem}
%%%%%%%%%%%%%%%%%%%%%%%%%%%%%%%%%%%%%%%%%%%%%%%%%%%%%%%%%%%%

\begin{theorem}[Fubini's Theorem]
Suppose

\[
f\in L^1(X\times Y,\mu\times\nu).
\]

Then the appropriate sections are integrable almost everywhere and

\[
\boxed{
\int_{X\times Y}
f\,d(\mu\times\nu)
=
\int_X
\left(
\int_Yf(x,y)\,d\nu(y)
\right)
d\mu(x)
}
\]

\[
\boxed{
=
\int_Y
\left(
\int_Xf(x,y)\,d\mu(x)
\right)
d\nu(y).
}
\]
\end{theorem}

Thus absolutely integrable functions may be integrated in either order.

%%%%%%%%%%%%%%%%%%%%%%%%%%%%%%%%%%%%%%%%%%%%%%%%%%%%%%%%%%%%
\subsection{Tonelli Versus Fubini}
\label{subsec:tonelli-versus-fubini}
%%%%%%%%%%%%%%%%%%%%%%%%%%%%%%%%%%%%%%%%%%%%%%%%%%%%%%%%%%%%

The distinction is important:

\[
\boxed{
\begin{array}{c|c}
\text{Tonelli} & f\geq0\\
\hline
\text{Fubini} & \int|f|<\infty
\end{array}
}
\]

Tonelli is often used first to determine whether

\[
\int|f|
\]

is finite.

If it is, Fubini then justifies changing the order of integration for the
signed function.

%%%%%%%%%%%%%%%%%%%%%%%%%%%%%%%%%%%%%%%%%%%%%%%%%%%%%%%%%%%%
\subsection{Worked Example: Fubini}
\label{subsec:worked-fubini}
%%%%%%%%%%%%%%%%%%%%%%%%%%%%%%%%%%%%%%%%%%%%%%%%%%%%%%%%%%%%

\begin{workedexamplebox}{Integral Over a Rectangle}

Let

\[
f(x,y)=xy
\]

on

\[
[0,1]\times[0,2].
\]

Since \(f\) is continuous on a compact rectangle, it is integrable.

Integrating first with respect to \(y\),

\[
\int_0^1
\left(
\int_0^2xy\,dy
\right)
dx.
\]

The inner integral is

\[
\int_0^2xy\,dy
=
x\left[\frac{y^2}{2}\right]_0^2
=
2x.
\]

Therefore

\[
\int_0^1 2x\,dx
=
1.
\]

\begin{answerbox}
\[
\boxed{
\int_{[0,1]\times[0,2]}
xy\,d(m\times m)
=
1.
}
\]
\end{answerbox}
\end{workedexamplebox}

%%%%%%%%%%%%%%%%%%%%%%%%%%%%%%%%%%%%%%%%%%%%%%%%%%%%%%%%%%%%
\subsection{Integral as Expectation}
\label{subsec:integral-as-expectation}
%%%%%%%%%%%%%%%%%%%%%%%%%%%%%%%%%%%%%%%%%%%%%%%%%%%%%%%%%%%%

On a probability space

\[
(\Omega,\mathcal{F},\mathbb{P}),
\]

the expectation of an integrable random variable \(X\) is simply its
Lebesgue integral:

\[
\boxed{
\mathbb{E}[X]
=
\int_\Omega X\,d\mathbb{P}.
}
\]

Thus expected value is not a separate mathematical operation.

It is integration with respect to probability measure.

%%%%%%%%%%%%%%%%%%%%%%%%%%%%%%%%%%%%%%%%%%%%%%%%%%%%%%%%%%%%
\subsection{Distribution Functions and Integration}
\label{subsec:distribution-integration}
%%%%%%%%%%%%%%%%%%%%%%%%%%%%%%%%%%%%%%%%%%%%%%%%%%%%%%%%%%%%

If a random variable has probability density \(p\) with respect to
Lebesgue measure, then

\[
\boxed{
\mathbb{P}(X\in A)
=
\int_Ap(x)\,dx.
}
\]

Its expectation is

\[
\boxed{
\mathbb{E}[X]
=
\int_{\R}x\,p(x)\,dx
}
\]

when the integral exists.

This demonstrates the direct connection between Lebesgue integration and
continuous probability distributions.

%%%%%%%%%%%%%%%%%%%%%%%%%%%%%%%%%%%%%%%%%%%%%%%%%%%%%%%%%%%%
\subsection{Signed Measures Preview}
\label{subsec:signed-measures-preview}
%%%%%%%%%%%%%%%%%%%%%%%%%%%%%%%%%%%%%%%%%%%%%%%%%%%%%%%%%%%%

A signed measure allows both positive and negative values while preserving
countable additivity.

A central structural result is the Jordan decomposition

\[
\boxed{
\nu=\nu^+-\nu^-,
}
\]

where

\[
\nu^+,
\qquad
\nu^-
\]

are positive measures.

This resembles the decomposition

\[
f=f^+-f^-
\]

for measurable functions.

%%%%%%%%%%%%%%%%%%%%%%%%%%%%%%%%%%%%%%%%%%%%%%%%%%%%%%%%%%%%
\subsection{Absolute Continuity of Measures}
\label{subsec:absolute-continuity-measures}
%%%%%%%%%%%%%%%%%%%%%%%%%%%%%%%%%%%%%%%%%%%%%%%%%%%%%%%%%%%%

Let \(\nu\) and \(\mu\) be measures.

We say that \(\nu\) is absolutely continuous with respect to \(\mu\) if

\[
\boxed{
\mu(A)=0
\Longrightarrow
\nu(A)=0.
}
\]

This is written

\[
\boxed{
\nu\ll\mu.
}
\]

The notation is read ``\(\nu\) is absolutely continuous with respect to
\(\mu\).''

%%%%%%%%%%%%%%%%%%%%%%%%%%%%%%%%%%%%%%%%%%%%%%%%%%%%%%%%%%%%
\subsection{Radon--Nikodym Theorem Preview}
\label{subsec:radon-nikodym-preview}
%%%%%%%%%%%%%%%%%%%%%%%%%%%%%%%%%%%%%%%%%%%%%%%%%%%%%%%%%%%%

One of the fundamental results of measure theory states that under
appropriate hypotheses, absolute continuity allows one measure to be
represented as an integral with respect to another.

\begin{theorem}[Radon--Nikodym Theorem]
Under standard sigma-finiteness assumptions, if

\[
\nu\ll\mu,
\]

then there exists a measurable function \(f\geq0\) such that

\[
\boxed{
\nu(A)
=
\int_Af\,d\mu
}
\]

for every measurable \(A\).

The function \(f\) is unique almost everywhere.
\end{theorem}

The function is denoted

\[
\boxed{
f=\frac{d\nu}{d\mu}
}
\]

and is called the Radon--Nikodym derivative.

This notation is read ``\(d\nu\) over \(d\mu\).''

%%%%%%%%%%%%%%%%%%%%%%%%%%%%%%%%%%%%%%%%%%%%%%%%%%%%%%%%%%%%
\subsection{Density as a Radon--Nikodym Derivative}
\label{subsec:density-radon-nikodym}
%%%%%%%%%%%%%%%%%%%%%%%%%%%%%%%%%%%%%%%%%%%%%%%%%%%%%%%%%%%%

Probability densities are examples of Radon--Nikodym derivatives.

If a probability measure \(\mathbb{P}\) on \(\R\) is absolutely continuous
with respect to Lebesgue measure \(m\), then

\[
\boxed{
p
=
\frac{d\mathbb{P}}{dm}
}
\]

is its density.

Consequently,

\[
\boxed{
\mathbb{P}(A)
=
\int_Ap\,dm.
}
\]

%%%%%%%%%%%%%%%%%%%%%%%%%%%%%%%%%%%%%%%%%%%%%%%%%%%%%%%%%%%%
\subsection{Common Errors}
\label{subsec:lebesgue-integration-common-errors}
%%%%%%%%%%%%%%%%%%%%%%%%%%%%%%%%%%%%%%%%%%%%%%%%%%%%%%%%%%%%

\subsubsection{Assuming Every Measurable Function Is Integrable}

Measurability alone does not imply

\[
\int|f|<\infty.
\]

For example,

\[
f(x)=1
\]

is measurable on \(\R\), but

\[
\int_{\R}1\,dm=\infty.
\]

%%%%%%%%%%%%%%%%%%%%%%%%%%%%%%%%%%%%%%%%%%%%%%%%%%%%%%%%%%%%
\subsubsection{Forgetting Absolute Integrability}

For a signed function, Lebesgue integrability means

\[
\boxed{
\int|f|\,d\mu<\infty.
}
\]

%%%%%%%%%%%%%%%%%%%%%%%%%%%%%%%%%%%%%%%%%%%%%%%%%%%%%%%%%%%%
\subsubsection{Treating \(\infty-\infty\) as Defined}

If

\[
\int f^+=\infty
\]

and

\[
\int f^-=\infty,
\]

then

\[
\int f
\]

is not defined by subtracting the two.

%%%%%%%%%%%%%%%%%%%%%%%%%%%%%%%%%%%%%%%%%%%%%%%%%%%%%%%%%%%%
\subsubsection{Ignoring Almost-Everywhere Equivalence}

Functions differing only on a null set have the same Lebesgue integral.

In \(L^p\), they represent the same equivalence class.

%%%%%%%%%%%%%%%%%%%%%%%%%%%%%%%%%%%%%%%%%%%%%%%%%%%%%%%%%%%%
\subsubsection{Using Monotone Convergence Without Monotonicity}

The Monotone Convergence Theorem requires a nonnegative increasing
sequence.

%%%%%%%%%%%%%%%%%%%%%%%%%%%%%%%%%%%%%%%%%%%%%%%%%%%%%%%%%%%%
\subsubsection{Using Dominated Convergence Without a Dominator}

Pointwise or almost-everywhere convergence is not enough.

One must find a single integrable function \(g\) satisfying

\[
|f_n|\leq g
\]

almost everywhere.

%%%%%%%%%%%%%%%%%%%%%%%%%%%%%%%%%%%%%%%%%%%%%%%%%%%%%%%%%%%%
\subsubsection{Confusing Fatou's Lemma with Equality}

Fatou's Lemma generally gives

\[
\boxed{
\int\liminf f_n
\leq
\liminf\int f_n,
}
\]

not equality.

%%%%%%%%%%%%%%%%%%%%%%%%%%%%%%%%%%%%%%%%%%%%%%%%%%%%%%%%%%%%
\subsubsection{Confusing \(L^p\) Convergence with Pointwise Convergence}

The condition

\[
\|f_n-f\|_p\to0
\]

does not mean

\[
f_n(x)\to f(x)
\]

at every point.

These are different notions of convergence.

%%%%%%%%%%%%%%%%%%%%%%%%%%%%%%%%%%%%%%%%%%%%%%%%%%%%%%%%%%%%
\subsubsection{Using Ordinary Supremum for \(L^\infty\)}

The \(L^\infty\) norm uses the essential supremum, not necessarily the
ordinary supremum.

%%%%%%%%%%%%%%%%%%%%%%%%%%%%%%%%%%%%%%%%%%%%%%%%%%%%%%%%%%%%
\subsubsection{Applying Fubini Without Integrability}

For signed functions, changing the order of integration generally requires

\[
\boxed{
\int|f|<\infty.
}
\]

For nonnegative functions, Tonelli's theorem is the appropriate tool.

%%%%%%%%%%%%%%%%%%%%%%%%%%%%%%%%%%%%%%%%%%%%%%%%%%%%%%%%%%%%
\subsection{Applications}
\label{subsec:lebesgue-integration-applications}
%%%%%%%%%%%%%%%%%%%%%%%%%%%%%%%%%%%%%%%%%%%%%%%%%%%%%%%%%%%%

\begin{applicationbox}
\textbf{Probability Theory.}
Expected values, moments, densities, and convergence of random variables
are formulated through Lebesgue integration.

\medskip

\textbf{Functional Analysis.}
The spaces \(L^p\) are fundamental Banach spaces, while \(L^2\) is one of
the central examples of a Hilbert space.

\medskip

\textbf{Fourier Analysis.}
Fourier transforms and Fourier series are naturally studied in \(L^1\)
and \(L^2\).

\medskip

\textbf{Harmonic Analysis.}
Integration, \(L^p\) estimates, maximal functions, and convergence
theorems form the basic language of harmonic analysis.

\medskip

\textbf{Partial Differential Equations.}
Weak solutions are frequently defined using integrals rather than
classical derivatives.

\medskip

\textbf{Sobolev Spaces.}
Weak derivatives are measured using \(L^p\) norms, producing the Sobolev
spaces used throughout modern PDE theory.

\medskip

\textbf{Statistics.}
Expectations, variances, likelihoods, and probability densities depend on
measure-theoretic integration.

\medskip

\textbf{Ergodic Theory.}
Long-term averages of dynamical systems are expressed using integrals with
respect to invariant measures.

\medskip

\textbf{Quantum Mechanics.}
Wave functions naturally belong to \(L^2\) spaces, where squared magnitude
is interpreted probabilistically.
\end{applicationbox}

%%%%%%%%%%%%%%%%%%%%%%%%%%%%%%%%%%%%%%%%%%%%%%%%%%%%%%%%%%%%
\subsection{Exercises}
\label{subsec:lebesgue-integration-exercises}
%%%%%%%%%%%%%%%%%%%%%%%%%%%%%%%%%%%%%%%%%%%%%%%%%%%%%%%%%%%%

\begin{exercise}[1]
Let \(E\) be measurable with

\[
\mu(E)=5.
\]

Compute

\[
\int_X3\mathbf{1}_E\,d\mu.
\]
\end{exercise}

\begin{exercise}[1]
Write

\[
f(x)
=
\begin{cases}
2, & x\in[0,1],\\
4, & x\in(1,3],\\
0, & \text{otherwise}
\end{cases}
\]

as a simple function using indicator functions.
\end{exercise}

\begin{exercise}[1]
Compute the Lebesgue integral of the function in the previous exercise.
\end{exercise}

\begin{exercise}[1]
Find \(f^+\) and \(f^-\) for

\[
f(x)=x
\]

on \([-1,1]\).
\end{exercise}

\begin{exercise}[1]
Explain why changing the value of a function at one point does not change
its Lebesgue integral on \(\R\).
\end{exercise}

\begin{exercise}[2]
Prove that if

\[
0\leq f\leq g,
\]

then

\[
\int f\,d\mu
\leq
\int g\,d\mu.
\]
\end{exercise}

\begin{exercise}[2]
Show that if \(f\geq0\) and

\[
\int f\,d\mu=0,
\]

then

\[
f=0
\]

almost everywhere.
\end{exercise}

\begin{exercise}[2]
Determine whether

\[
f(x)=\frac1x
\]

belongs to \(L^1([1,\infty))\).
\end{exercise}

\begin{exercise}[2]
Determine whether

\[
f(x)=\frac1{x^2}
\]

belongs to \(L^1([1,\infty))\).
\end{exercise}

\begin{exercise}[2]
Determine for which \(p\) the function

\[
f(x)=x^{-1/2}
\]

belongs to \(L^p((0,1))\).
\end{exercise}

\begin{exercise}[2]
For

\[
f_n(x)=x^n
\]

on \([0,1]\), determine the pointwise limit and use the Dominated
Convergence Theorem to compute

\[
\lim_{n\to\infty}
\int_0^1x^n\,dx.
\]
\end{exercise}

\begin{exercise}[2]
Use the Monotone Convergence Theorem to justify

\[
\int_X
\sum_{n=1}^{\infty}f_n\,d\mu
=
\sum_{n=1}^{\infty}
\int_Xf_n\,d\mu
\]

for nonnegative measurable \(f_n\).
\end{exercise}

\begin{exercise}[2]
Show that the Dirichlet function

\[
\mathbf{1}_{\Q}
\]

has Lebesgue integral zero on \([0,1]\).
\end{exercise}

\begin{exercise}[3]
Prove Fatou's Lemma using the Monotone Convergence Theorem.
\end{exercise}

\begin{exercise}[3]
Derive the Dominated Convergence Theorem from Fatou's Lemma.
\end{exercise}

\begin{exercise}[3]
Let

\[
f_n(x)
=
n\mathbf{1}_{(0,1/n)}(x).
\]

Show that \(f_n\to0\) almost everywhere but not in \(L^1(0,1)\).
\end{exercise}

\begin{exercise}[3]
Show that if

\[
f_n\to f
\]

in \(L^1\), then

\[
\int f_n\,d\mu
\to
\int f\,d\mu.
\]
\end{exercise}

\begin{exercise}[3]
Prove that \(L^p\) convergence implies convergence in measure for
\(1\leq p<\infty\).
\end{exercise}

\begin{exercise}[3]
Show that if \(f\in L^\infty\) and \(g\in L^1\), then

\[
fg\in L^1
\]

and

\[
\boxed{
\|fg\|_1
\leq
\|f\|_\infty\|g\|_1.
}
\]
\end{exercise}

\begin{exercise}[3]
Use H\"older's inequality to prove that if \(\mu(X)<\infty\) and
\(1\leq p<q\leq\infty\), then

\[
L^q(X)\subseteq L^p(X).
\]
\end{exercise}

\begin{exercise}[3]
Use Fubini's theorem to evaluate

\[
\int_0^1\int_0^2xy\,dy\,dx.
\]
\end{exercise}

\begin{exercise}[3]
Use Tonelli's theorem to justify interchanging the order of summation and
integration for a series of nonnegative measurable functions.
\end{exercise}

\begin{exercise}[3]
Prove the absolute continuity of the Lebesgue integral for
\(f\in L^1\).
\end{exercise}

\begin{exercise}[4]
Prove H\"older's inequality.
\end{exercise}

\begin{exercise}[4]
Use H\"older's inequality to prove Minkowski's inequality.
\end{exercise}

\begin{exercise}[4]
Prove that the \(L^p\) norm is well defined on equivalence classes modulo
equality almost everywhere.
\end{exercise}

\begin{exercise}[4]
Study the proof that \(L^p\) is complete for

\[
1\leq p\leq\infty.
\]
\end{exercise}

\begin{exercise}[4]
Construct a sequence that converges in measure but does not converge almost
everywhere.
\end{exercise}

\begin{exercise}[4]
Give an example of almost-everywhere convergence that does not imply
\(L^1\) convergence.
\end{exercise}

\begin{exercise}[4]
Investigate the hypotheses of Fubini's theorem and construct an example
showing why absolute integrability is important.
\end{exercise}

\begin{exercise}[4]
Study the Radon--Nikodym theorem and explain how probability density
functions arise as Radon--Nikodym derivatives.
\end{exercise}

%%%%%%%%%%%%%%%%%%%%%%%%%%%%%%%%%%%%%%%%%%%%%%%%%%%%%%%%%%%%
\subsection{Conceptual Summary}
\label{subsec:lebesgue-integration-summary}
%%%%%%%%%%%%%%%%%%%%%%%%%%%%%%%%%%%%%%%%%%%%%%%%%%%%%%%%%%%%

Lebesgue integration builds integration directly from measure.

For measurable \(E\),

\[
\boxed{
\int\mathbf{1}_E\,d\mu
=
\mu(E).
}
\]

For a nonnegative simple function,

\[
\phi
=
\sum_{k=1}^{n}
a_k\mathbf{1}_{E_k},
\]

the integral is

\[
\boxed{
\int\phi\,d\mu
=
\sum_{k=1}^{n}
a_k\mu(E_k).
}
\]

For a nonnegative measurable function,

\[
\boxed{
\int f\,d\mu
=
\sup_{\substack{0\leq\phi\leq f\\
\phi\text{ simple}}}
\int\phi\,d\mu.
}
\]

A signed function is decomposed as

\[
\boxed{
f=f^+-f^-.
}
\]

Lebesgue integrability is characterized by

\[
\boxed{
\int|f|\,d\mu<\infty.
}
\]

The Monotone Convergence Theorem states

\[
\boxed{
0\leq f_n\uparrow f
\Longrightarrow
\int f_n\,d\mu
\to
\int f\,d\mu.
}
\]

Fatou's Lemma gives

\[
\boxed{
\int\liminf f_n\,d\mu
\leq
\liminf\int f_n\,d\mu.
}
\]

The Dominated Convergence Theorem states that if

\[
f_n\to f
\quad\text{a.e.}
\]

and

\[
|f_n|\leq g\in L^1,
\]

then

\[
\boxed{
\int f_n\,d\mu
\to
\int f\,d\mu.
}
\]

The \(L^p\) spaces are defined through

\[
\boxed{
\|f\|_p
=
\left(
\int|f|^p\,d\mu
\right)^{1/p}
}
\]

for

\[
1\leq p<\infty,
\]

while

\[
\boxed{
\|f\|_\infty
=
\operatorname*{ess\,sup}|f|.
}
\]

H\"older's inequality controls products:

\[
\boxed{
\|fg\|_1
\leq
\|f\|_p\|g\|_q.
}
\]

Minkowski's inequality provides the triangle inequality:

\[
\boxed{
\|f+g\|_p
\leq
\|f\|_p+\|g\|_p.
}
\]

Product measures permit multidimensional integration.

Tonelli's theorem handles nonnegative measurable functions, while Fubini's
theorem handles absolutely integrable functions.

Lebesgue integration therefore provides a flexible framework in which
integration, infinite limits, probability, and function spaces interact
naturally.

%%%%%%%%%%%%%%%%%%%%%%%%%%%%%%%%%%%%%%%%%%%%%%%%%%%%%%%%%%%%
\subsection{Section Connections}
\label{subsec:lebesgue-integration-connections}
%%%%%%%%%%%%%%%%%%%%%%%%%%%%%%%%%%%%%%%%%%%%%%%%%%%%%%%%%%%%

\chapterconnection{
Lebesgue Integration builds directly on Measure Theory by using measurable
sets and measurable functions to define a substantially more flexible
integral than the Riemann integral. The Monotone Convergence Theorem,
Fatou's Lemma, and the Dominated Convergence Theorem provide powerful
conditions for interchanging limits and integrals. The resulting \(L^p\)
spaces connect integration to geometry and Functional Analysis, with
\(L^2\) becoming a central Hilbert space. Product measures and the
Tonelli--Fubini theorems provide the rigorous foundation for multiple
integration. Probability Theory interprets expectation as Lebesgue
integration with respect to probability measure, while Fourier Analysis,
Harmonic Analysis, Partial Differential Equations, Sobolev spaces, and
Operator Theory rely extensively on \(L^p\) methods. The next section,
Complex Analysis, develops analysis over the complex numbers, while later
sections return to \(L^p\) spaces from the perspectives of Functional and
Harmonic Analysis.
}

%%%%%%%%%%%%%%%%%%%%%%%%%%%%%%%%%%%%%%%%%%%%%%%%%%%%%%%%%%%%
% SECTION SOURCE TRACKING
%%%%%%%%%%%%%%%%%%%%%%%%%%%%%%%%%%%%%%%%%%%%%%%%%%%%%%%%%%%%

%------------------------------------------------------------
% 5.7 Lebesgue Integration
%------------------------------------------------------------
%
% Source basis:
%   - Standard undergraduate and beginning graduate real analysis
%     and measure-theoretic integration.
%   - Standard development of the Lebesgue integral and L^p spaces.
%
% Used for:
%   - indicator functions;
%   - simple functions;
%   - simple-function approximation;
%   - nonnegative measurable integrals;
%   - positive and negative parts;
%   - integrable functions;
%   - almost-everywhere equality;
%   - Monotone Convergence Theorem;
%   - Fatou's Lemma;
%   - Dominated Convergence Theorem;
%   - Riemann versus Lebesgue integration;
%   - Lebesgue criterion for Riemann integrability;
%   - absolute continuity of the integral;
%   - L^1 and L^p spaces;
%   - essential supremum and L^\infty;
%   - Holder's inequality;
%   - Minkowski's inequality;
%   - completeness of L^p;
%   - L^p convergence;
%   - product measures;
%   - Tonelli's theorem;
%   - Fubini's theorem;
%   - expectation as integration;
%   - signed measures preview;
%   - absolute continuity of measures;
%   - Radon--Nikodym theorem preview.
%
% Notes:
%   - Sigma-algebras, measures, measurable functions, null sets,
%     and Lebesgue measure are developed in Section 5.6.
%   - Riemann integration is developed in Sections 5.2 and 5.5.
%   - Probability measures and random variables are developed
%     extensively in Chapter 7.
%   - Banach and Hilbert space structures of L^p spaces are
%     developed further in Functional Analysis.
%   - Fourier and Harmonic Analysis make extensive later use of
%     L^1 and L^2 methods.
%
%%%%%%%%%%%%%%%%%%%%%%%%%%%%%%%%%%%%%%%%%%%%%%%%%%%%%%%%%%%%